\subsection{Intervalle}
\label{subsec:fundamentals_intervals}

\begin{DefinitionBox}{Intervall}{fundamentals_interval}
  Ein \emph{Intervall} \EnglishTerm{interval} ist eine zusammenhängende
  Teilmenge $I\subseteq\R$, die zwischen zwei Randpunkten $a<b$ liegt und alle
  reellen Zahlen dazwischen enthält.

  \begin{center}
    \CachedTikzFigure{content/resources/chapters/fundamentals/tikz/interval_definition_number_line.tex}
  \end{center}
\end{DefinitionBox}

Reelle Randpunkte können zum Intervall gehören oder ausgeschlossen sein:\par

\renewcommand{\arraystretch}{1.35}
\begin{tabularx}{\linewidth}{@{}l >{\centering\arraybackslash}X l@{}}
  \toprule
  \textbf{Intervalltyp} & \textbf{Mengenschreibweise} & \textbf{Randpunkte} \\
  \midrule
  abgeschlossen
    & $[a,b]=\{x\in\R\mid a\,\Cloze{0.5cm}{\leq}\,x\,\Cloze{0.5cm}{\leq}\,b\}$
    & $a,b$ enthalten \\
  offen
    & $(a,b)=\{x\in\R\mid a\,\Cloze{0.5cm}{<}\,x\,\Cloze{0.5cm}{<}\,b\}$
    & $a,b$ nicht enthalten \\
  links abgeschlossen
    & $[a,b)=\{x\in\R\mid a\,\Cloze{0.5cm}{\leq}\,x\,\Cloze{0.5cm}{<}\,b\}$
    & nur $a$ enthalten \\
  rechts abgeschlossen
    & $(a,b]=\{x\in\R\mid a\,\Cloze{0.5cm}{<}\,x\,\Cloze{0.5cm}{\leq}\,b\}$
    & nur $b$ enthalten \\
  \bottomrule
\end{tabularx}
\renewcommand{\arraystretch}{1}

\begin{RemarkBox}[remark title=Alternative Intervallschreibweise]
  Für offene und halboffene Intervalle ist im deutschsprachigen Raum auch
  folgende Schreibweise üblich:
  \[
    (a,b)=]a,b[,
    \qquad
    [a,b)=[a,b[,
    \qquad
    (a,b]=]a,b].
  \]
\end{RemarkBox}


\begin{KeyPointBox}{Uneigentliche Randpunkte}{fundamentals_infinite_interval_endpoints}
  Liegt ein Randpunkt im Unendlichen, spricht man von einem uneigentlichen Randpunkt.
  Da $\pm \infty$ keine reellen Zahlen sind, gehören sie daher nie zum Intervall.
  An einem uneigentlichen Randpunkt (im Unendlichen) ist ein Intervall also stets offen.

  \begin{center}
    \renewcommand{\arraystretch}{1.45}
    \begin{tabularx}{0.86\linewidth}{@{}>{\centering\arraybackslash}X >{\centering\arraybackslash}X@{}}
      $(-\infty,b]$\quad\textcolor{green!50!black}{\textbf{richtig}}
      & $[a,\infty)$\quad\textcolor{green!50!black}{\textbf{richtig}} \\
      $[-\infty,b]$\quad\textcolor{red!75!black}{\textbf{falsch}}
      & $(a,\infty]$\quad\textcolor{red!75!black}{\textbf{falsch}}
    \end{tabularx}
    \renewcommand{\arraystretch}{1}
  \end{center}
\end{KeyPointBox}
